\documentclass[11pt,a4paper]{article}

\usepackage[utf8]{inputenc}
\usepackage[T1]{fontenc}
\usepackage{lmodern}
\usepackage{amsmath,amssymb}
\usepackage{booktabs,tabularx,graphicx}
\usepackage{float}
\usepackage{geometry,microtype}
\usepackage{xcolor}
\usepackage{fancyhdr}

\geometry{left=2.3cm,right=2.3cm,top=2.3cm,bottom=2.2cm}
\graphicspath{{figuras/}}
\definecolor{heading}{HTML}{17324D}
\setlength{\parindent}{0pt}
\setlength{\parskip}{0.55em}
\renewcommand{\arraystretch}{1.2}
\renewcommand{\figurename}{Figura}
\renewcommand{\tablename}{Tabela}
\pagestyle{fancy}
\fancyhf{}
\fancyhead[L]{\small Questão 4.3 --- Bike Sharing}
\fancyfoot[C]{\small\thepage}
\renewcommand{\headrulewidth}{0.4pt}

\begin{document}

{\Large\bfseries\color{heading}Questão 4.3 --- Temperatura e baixa utilização}\par

\section*{Critério e preparação dos dados}

A análise mantém a seleção definida na Questão 1: com $M=581706$,
$r=1+(M\bmod 100)=7$, sendo utilizadas as $300$ observações consecutivas
de 07/01/2011 a 02/11/2011. O total diário e a classificação seguem as
definições da Questão 2:
\[
U_i=\texttt{casual}_i+\texttt{registered}_i,
\qquad Q_1(U)=2105{,}5,
\]
\[
\texttt{low\_usage}_i=
\begin{cases}
1, & U_i<Q_1(U),\\
0, & U_i\geq Q_1(U).
\end{cases}
\]
O quartil foi calculado sobre os $300$ dias pelo método
\texttt{quantile(type = 7)} do R, antes de dividir as observações em grupos.
Assim, $75$ dias ($25\%$) são classificados como de baixa utilização e os
$225$ restantes pertencem ao outro grupo. Não se recalcula um quartil para
cada grupo ou estação.

Para tornar o eixo horizontal interpretável, a temperatura normalizada do
enunciado é convertida por $T_{\mathrm{C}}=41T_{\mathrm{norm}}$. O script
também reconhece uma coluna que já esteja em graus Celsius. Nenhuma
observação é excluída ou reordenada.

\section*{Dispersão e classificação dos dias}

A Figura~\ref{fig:dispersao} representa a temperatura e o total de usuários
em cada dia. A linha horizontal marca o mesmo $Q_1$ usado na classificação:
os pontos abaixo dela pertencem a \texttt{low\_usage = 1}. A concentração
dos dias de baixa utilização ocorre predominantemente em temperaturas
menores, enquanto os totais mais elevados são observados com maior
frequência em temperaturas moderadas a altas.

\begin{figure}[H]
\centering
\includegraphics[width=0.84\linewidth]{questao_4_3_dispersao.png}
\caption{Temperatura e total diário, com os $300$ dias diferenciados pela
classificação de baixa utilização. A linha tracejada é $Q_1=2105{,}5$.
Gráfico produzido pelo script \texttt{questao\_4\_3.R}.}
\label{fig:dispersao}
\end{figure}

\clearpage
\section*{Comparação visual entre os grupos}

A Figura~\ref{fig:paineis} separa os dois grupos, preservando as mesmas
escalas nos eixos. As retas descrevem a tendência linear \emph{dentro de
cada subconjunto} e foram desenhadas apenas no intervalo de temperatura
observado no respectivo grupo. A comparação mostra uma tendência positiva
aparentemente menos marcada entre os dias de baixa utilização.

\begin{figure}[H]
\centering
\includegraphics[width=0.80\linewidth]{questao_4_3_paineis.png}
\caption{Temperatura e utilização nos dois grupos. As escalas são iguais;
as retas representam ajustes lineares descritivos. Gráfico produzido pelo
script \texttt{questao\_4\_3.R} a partir da tabela do PDF, cuja temperatura
é arredondada. Execute o script com o CSV original para os valores finais.}
\label{fig:paineis}
\end{figure}

\begin{table}[H]
\centering
\small
\caption{Temperatura e correlação de Pearson em cada grupo.}
\begin{tabular}{@{}lrrrr@{}}
\toprule
Grupo & Dias & Média de $T$ & Intervalo de $T$ & $r(T,U)$ \\
\midrule
Baixa utilização & 75 & $11{,}08\,^{\circ}\mathrm C$ & $2{,}4$--$27{,}9\,^{\circ}\mathrm C$ & $0{,}277$ \\
Demais dias & 225 & $24{,}50\,^{\circ}\mathrm C$ & $10{,}8$--$34{,}8\,^{\circ}\mathrm C$ & $0{,}568$ \\
\bottomrule
\end{tabular}
\end{table}

\section*{Interpretação e limitações}

O coeficiente de Pearson é positivo em ambos os grupos, mas menor entre os
dias classificados como de baixa utilização ($r\approx0{,}277$) do que nos
demais ($r\approx0{,}568$). As temperaturas médias também diferem:
$11{,}08\,^{\circ}\mathrm C$ e $24{,}50\,^{\circ}\mathrm C$,
respectivamente. A presença de dias quentes com pouco uso impede uma leitura
determinista da relação: em 27/08/2011, por exemplo, constam aproximadamente
$27{,}9\,^{\circ}\mathrm C$ e $1115$ usuários.

Essa comparação é descritiva. A variável \texttt{low\_usage} foi construída
a partir do próprio total de usuários, representado no eixo vertical. A
separação restringe por definição a faixa de $U$ em cada grupo; diferenças
entre correlações ou inclinações calculadas depois do corte não demonstram
efeitos causais distintos da temperatura. Estação do ano e condições
meteorológicas também podem estar associadas à utilização. Na amostra,
$60$ dos $75$ dias de baixa utilização ocorreram no inverno.

Nas dez primeiras observações do grupo, todos os dias satisfazem
$U_i<2105{,}5$. Portanto, não existe nesse recorte inicial uma segunda
classe com a qual comparar a relação. A avaliação entre os dois grupos
depende das $300$ observações completas.

\end{document}
